% titel: Blatt 5, A6: Vergleich des Ionenradius

\subsection*{Frage}
Welches Ion hat je Paar den grösseren Ionenradius und warum?

\subsection*{Lösung}
\begin{itemize}
  \item \textbf{$\mathrm{K^+}$, $\mathrm{Sc^{3+}}$:} Beide haben 18 Elektronen
    (isoelektronisch zu $[\mathrm{Ar}]$). $\mathrm{K^+}$ ($Z=19$) hat die
    geringere Kernladung als $\mathrm{Sc^{3+}}$ ($Z=21$) $\to$ $\mathrm{K^+}$ ist
    grösser.
  \item \textbf{$\mathrm{Cu^+}$, $\mathrm{Ag^+}$:} Gleiche Ladung, Ag liegt eine
    Periode tiefer $\to$ mehr Schalen $\to$ $\mathrm{Ag^+}$ ist grösser.
  \item \textbf{$\mathrm{Se^{2-}}$, $\mathrm{Br^-}$:} Beide 36 Elektronen
    (isoelektronisch zu $[\mathrm{Kr}]$). $\mathrm{Se^{2-}}$ ($Z=34$) hat die
    geringere Kernladung als $\mathrm{Br^-}$ ($Z=35$) $\to$ $\mathrm{Se^{2-}}$
    ist grösser.
  \item \textbf{$\mathrm{O^{2-}}$, $\mathrm{S^{2-}}$:} Gleiche Ladung, S liegt
    eine Periode tiefer $\to$ mehr Schalen $\to$ $\mathrm{S^{2-}}$ ist grösser.
  \item \textbf{$\mathrm{Ca^{2+}}$, $\mathrm{P^{3-}}$:} Beide 18 Elektronen
    (isoelektronisch). $\mathrm{P^{3-}}$ ($Z=15$) hat die geringste Kernladung,
    $\mathrm{Ca^{2+}}$ ($Z=20$) die grösste $\to$ $\mathrm{P^{3-}}$ ist am
    grössten.
\end{itemize}

\begin{ergebnis}
grösser: $\mathrm{K^+}$ (über $\mathrm{Sc^{3+}}$), $\mathrm{Ag^+}$ (über
$\mathrm{Cu^+}$), $\mathrm{Se^{2-}}$ (über $\mathrm{Br^-}$), $\mathrm{S^{2-}}$
(über $\mathrm{O^{2-}}$), $\mathrm{P^{3-}}$ (über $\mathrm{Ca^{2+}}$).
\end{ergebnis}

\begin{kontrolle}
Alle Ergebnisse folgen konsistent aus zwei Regeln: isoelektronisch $\to$
kleinere Kernladung grösser; gleiche Ladung $\to$ tiefere Periode grösser.
Beispiel $\mathrm{O^{2-}}$ (gross, Anion) gegen $\mathrm{Ca^{2+}}$ (klein,
Kation, hohe Ladung) ist plausibel für die Grössenordnung Ionenkristallgitter.
\end{kontrolle}

\begin{merksatz}
Bei gleicher Elektronenzahl entscheidet die Kernladung: weniger $Z$ = grösser
(= negativeres Ion). Bei gleicher Ladung entscheidet die Periode: tiefer =
grösser.
\end{merksatz}
